\paragraph{How many dimensions does the chaos have?} For chaotic attractors the corresponding measure is the
Kaplan--Yorke dimension \cite{kaplan1979chaotic}, $D_{\mathrm{KY}}=j+\sum_{i\le j}\lambda_i/|\lambda_{j+1}|$ with $j$ the
largest index at which the partial sum of the ordered Lyapunov exponents is still nonnegative; with one positive
exponent and $\lambda_1+\lambda_2<0$, $D_{\mathrm{KY}}=1+\lambda_1/|\lambda_2|$. We computed the four leading exponents
by QR iteration (Section~\ref{sec:methods}; 1000 transient and 2000 measured steps, float64) at chaotic cells chosen
by percentile of their largest exponent, before any of their spectra were seen. Every attractor we measured has
exactly one positive exponent, so none is hyperchaotic, but the dimension grows with temperature. Along the hostile
temperature line, where we fit a scalar return relation (Section~\ref{sec:hostility}), eight chaotic cells have
$D_{\mathrm{KY}}=1.008$ to $1.16$: one-dimensional to within a sixth of a dimension. At fourteen chaotic cells of the
headline maps ($\log_{10}T$ from $-0.17$ to $-0.04$) it is $1.04$ to $1.65$, rising with $\lambda_1$. In the
high-temperature region ($\log_{10}T$ from $-0.02$ to $0.13$), nine chaotic cells have $D_{\mathrm{KY}}=1.37$ to
$2.02$; at three of them $\lambda_1+\lambda_2\ge0$, and the attractor is effectively two-dimensional. A scalar return
map describes the low-temperature chaos well and the high-temperature chaos poorly, consistent with the large second
multipliers of the high-temperature bubble.
\input{figures/tikz-atlas}

\section{Methods}
\label{sec:methods}